% /root/monodromy, 2026-11-19. Root, reduction-descent and obstructions-crystals complete independent proof/source/code audits passed.
\subsection{All degrees for a generic mixed elliptic product}
\label{subsec:generic-mixed-product-groups-candidate}
\status{Proved and independently checked. The parent, reduction/descent and obstructions/crystals agents audited the complete proof and independently reran the coefficient reconstruction. The ordinary elliptic modulus is assumed transcendental over the prime field; every fixed algebraic ordinary modulus remains outside this theorem.}

Let $E_0:y^2+y=x^3$ and let $E$ be an ordinary elliptic curve
in characteristic two. Put
\[
 X=E\times E_0,\qquad
 L=\mathcal O_E(3O)\boxtimes\mathcal O_{E_0}(3O),\qquad
 b_d=9d^2.
\]
The rank and entire invariant-covector results already proved give
relative ranks $b_d-1$ in odd degree and $b_d-2$ in even degree,
with respectively zero and two geometric invariant covectors.
Here and below a sharp curve means a smooth curve attaining the
indicated maximal rank; on a mixed product its full curve defect
$g-f$ is one in odd degree and two in even degree.

\begin{lemma}[Generic odd images and odd Hasse divisors]
\label{lem:generic-mixed-odd-groups-candidate}
If $j(E)$ is transcendental over $\mathbf F_2$, the geometric
relative degree-$d$ image is $\GL_{b_d-1}(\mathbf F_2)$ for every
odd $d\ge1$. For odd $d\ge3$, the determinant on the relative
coherent bundle is a geometrically irreducible reduced homogeneous
polynomial of degree $b_d$, with generic stable loss one on its
divisor.
\end{lemma}

\begin{proof}
Consider the one-parameter family
\[
 E_t:\quad y^2+txy+(1+t)y=x^3+1,
 \qquad \Delta(t)=1+t^3+t^5,
 \qquad j(t)=\frac{t^{12}}{1+t^3+t^5}.
\]
Work over $B=\Spec\mathbf F_2[t,\Delta(t)^{-1}]$.
Its fibres at $t=0,1$ are smooth; the first is geometrically
isomorphic to $E_0$ and the second is
$y^2+xy=x^3+1$. The displayed nonconstant $j$ makes the generic
ordinary modulus transcendental. The standard pole-order bases
of $H^0(E_t,\mathcal O(3dO))$ trivialize its coefficient projective
bundle over $B$.

For odd $d$, the relative coherent bundle on each coefficient
projective space is $\Omega^1$, and it is generically Frobenius
bijective by Theorem~\ref{thm:DE-A3-ppav-ranks}. This construction
and its rank-$b_d-1$ Lang frame cover extend over the open of
Frobenius invertibility in the relative projective bundle over
$B$. The geometric $t=0$ frame fibre is connected, by
Theorem~\ref{thm:ss-all-odd-full-groups}. The geometric generic
frame fibre is therefore connected as well. Here no properness
of the invertibility open is assumed: normalize $B$ in a finite
separable extension defining the hypothetical generic components.
The normalization is finite, since $B$ is excellent. Its pulled-back
coefficient open is normal and integral. Components of a finite
etale cover of this open are disjoint, open and closed, and each
surjects onto the open. All would consequently meet a geometric
fibre over $t=0$, contradicting that fibre's connectedness.
This is precisely the nonproper argument of
Lemma~\ref{lem:ordinary-pair-GL8-open}.
Connectedness of the full frame cover proves the asserted generic
general linear group.

The determinant on $\Omega^1$ has degree $b_d$ and varies
algebraically with $t$. For $d\ge3$ its $t=0$ polynomial is
geometrically irreducible and reduced, by
Theorem~\ref{thm:ss-all-odd-Hasse}. Geometric reducibility of
a nonzero homogeneous form of fixed degree is a closed condition:
for each possible pair of positive factor degrees it is the image
of the proper multiplication map between the corresponding
projective spaces of forms. Their finite union includes every
nonreduced form. Thus the geometric generic polynomial is
irreducible and reduced. A point of stable loss one on the
special divisor belongs to the open where an appropriate
$(b_d-2)$-minor of a sufficiently high Frobenius iterate is
nonzero. The total determinant divisor is flat over a
neighbourhood of that point: its defining form is nonzero on
every fibre. Therefore this nonempty open meets the generic
divisor, giving generic stable loss one. Equivalently one may
use the first nonzero minor and the absence of a vertical
component of this flat hypersurface.

Finally any elliptic curve with transcendental $j$ becomes
isomorphic over an algebraic closure to a base change of this
geometric generic curve, since $j(t)$ is nonconstant.
Finite etale connectedness is preserved under algebraically
closed field extension. This proves each assertion for that
curve. The same generic point works for all $d$; no intersection
of countably many closed-point open conditions is being taken.
\end{proof}

\begin{lemma}[The mixed first-image sheaf and a natural block]
\label{lem:generic-mixed-image-block-candidate}
Work over an algebraically closed field, and suppose the odd
degree-one image on $X$ is $\GL_8(\mathbf F_2)$.
For even $d\ge2$ put $N=L^d$, $W=H^0(N)$, $b=b_d$,
$P=\mathbf P(W)$ and $\mathcal B=W^*\otimes\mathcal O_P(-1)$.
Normalize invariant derivations $e,n$ so that
$e^2=e$, $n^2=0$ and $en=ne$, with $e$ from $E$ and $n$ from
$E_0$. They lift canonically to $N$, since $N$ is a Frobenius
pullback. Set
\[
 \mathcal I=\ker\bigl(\mathcal B\xrightarrow{(s,es)}\mathcal O_P^2\bigr),
 \qquad
 \mathcal K=\ker\bigl(\mathcal B\xrightarrow{(s,es,ns,ens)}\mathcal O_P^4\bigr).
\]
There is a closed subset $Z\subset P$, of codimension $b/4$,
outside which these are vector bundles and
\[
 0\longrightarrow\mathcal K\longrightarrow\mathcal I
       \xrightarrow{(ns,ens)}\mathcal O_{P\setminus Z}^2
       \longrightarrow0
\]
is Frobenius-equivariant with constant unit quotient. Their
reflexive extensions satisfy
\[
 \rank\mathcal I=b-2,\quad\rank\mathcal K=b-4,\quad
 \det\mathcal I=\det\mathcal K=\mathcal O_P(-b),\quad
 H^0(P,\mathcal I)=H^0(P,\mathcal K)=0.
\]
Frobenius on $\mathcal I$ is generically bijective, its Lang
system is the required relative unit system, and its determinant
is a nonzero degree-$b$ form $h_d^{\mathrm u}$.
The geometric group contains an actual natural
$\GL_8(\mathbf F_2)$ block fixing a complementary subspace.
Its natural eight-space lies in the common covector kernel.
\end{lemma}

\begin{proof}
On the ordinary factor,
$F_*\mathcal O_E=\mathcal O\oplus B_E$ with $B_E$ the nontrivial
geometric two-torsion line bundle. The two summands are the
weight-zero and weight-one spaces for $e$. On the supersingular
factor the positive twist of the nonsplit extension
$0\to\mathcal O\to F_*\mathcal O_{E_0}\to\mathcal O\to0$
has vanishing first cohomology; hence its global sections are
free over $k[n]/n^2$. Consequently
\[
 W=W_0\oplus W_1,\qquad e|W_i=i,
\]
and each $W_i$ is free of rank $b/4$ over $k[n]/n^2$.
Write $s=s_0+s_1$. The locus
\[
 Z=\{ns_0=0\}\cup\{ns_1=0\}
\]
is the union of two linear subspaces of codimension $b/4$.
Outside it, $s_0,ns_0$ and $s_1,ns_1$ are independent in their
respective weight spaces. This proves independence of
$s,es,ns,ens$ and all the claimed bundle ranks and surjectivity.
The kernel sheaves are reflexive by the depth lemma, and
codimension at least two permits their determinant calculation
from the displayed maps. Their inclusions into $\mathcal B$
give the vanishing of global sections.

The mixed row of \eqref{eq:nonordinary-Cartier-table} says
\[
 M(s)s=M(s)es=0,\qquad
 M(s)ns=(ns)^{[2]},\qquad M(s)ens=(ens)^{[2]},
\]
in the root-removed Cartier convention. Thus full Frobenius
maps into $\mathcal I$ and the two quotient covectors are
Frobenius-fixed. The exact generic relative stable rank is
$b-2$, so its restriction to $\mathcal I$ is bijective on a
dense open. The full coherent model has the same unit part
as the relative system: the ambient top coherent cohomology
is Frobenius-nilpotent. Its determinant extends across $Z$
as a nonzero section of $\mathcal O_P(b)$, and agrees with
the determinant on $\mathcal K$ because the quotient is unit.

For the block, fix a smooth sharp odd-degree curve
$D\in|L^{d-1}|$ and vary a sharp $C\in|L|$ meeting it
transversely. The full curve defect of their nodal union is
$1+1=2$, precisely the generic even defect. At an ordinary
intersection node of its equation $s$, the first derivatives
vanish and the mixed derivative $ens$ is nonzero. Thus both
$ns_1=ens$ and $ns_0=ns+ens$ are nonzero there, placing the
union outside $Z$. The boundary Lang system therefore embeds
with the full generic rank, as in
Theorem~\ref{thm:ss-uniform-GL8-block}.

There are $q=18(d-1)$ intersection points. Their monodromy
is $S_q$: the one- and two-point incidences on the smooth
integral curve $D$ are projective bundles, using respectively
one and two independent values. The double-zero incidence
has a general unique double zero with independent value and
first-jet deformation parameters; its quadratic local algebra
$z^2+b'z+c$ supplies a single transposition, including in
characteristic two. The joint curve-jet rank computation for
two plane cubics excludes a second double zero, exactly as in
Lemma~\ref{lem:ss-intersection-Hasse-degree}. Its hypotheses
also hold for an ordinary plane cubic: its quadratic partials
are independent, so the tangents through a fixed point form a
finite set. Removing the sharp-source complement does not
change the image of a finite cover of a normal integral base.

The normalization filtration of the boundary relative unit
system $V$ has graph term $T$ and a quotient
\[
 V/T=(H^1_{\et}(C,\mathbf F_2)\oplus
          H^1_{\et}(D,\mathbf F_2))/H^1_{\et}(X,\mathbf F_2).
\]
The ambient diagonal line maps injectively to either curve.
Projection to the relative degree-one system on $C$ has
kernel an extension $T_0$ of the graph term by a constant
space. The image on $T_0$ has a normal two-group with quotient
a subgroup of $S_q$; on its faithful graph quotient it is the
full $S_q$. Hence its only possible nonabelian simple factor
is $A_q$. It has no factor $\GL_8(\mathbf F_2)$: the latter
is simple, has no prime factor eleven, whereas $A_q$ does
since $q\ge18$.

Goursat's lemma therefore gives a subgroup still surjecting
onto $\GL_8$ and acting trivially on $T_0$. In the resulting
block parabolic its translation kernel is a sum of natural-dual
copies. The elementary vanishing
$H^1(\GL_8(\mathbf F_2),(\mathbf F_2^8)^*)=0$
splits its quotient complement after conjugation; removing
the remaining translation columns gives an actual natural
$\GL_8$ block fixing a complement, as proved in the cited
block theorem. The two invariant covectors vanish on this
nontrivial simple eight-dimensional block, so it lies in their
common kernel.
\end{proof}

\begin{lemma}[The mixed even Hasse divisor in degrees at least four]
\label{lem:generic-mixed-even-Hasse-candidate}
Suppose $j(E)$ is transcendental over $\mathbf F_2$.
For every even $d\ge4$, the polynomial $h_d^{\mathrm u}$ of
the preceding lemma is geometrically irreducible and reduced,
and its generic divisor has stable loss exactly one.
\end{lemma}

\begin{proof}
We spell out the changes to the two boundary degree comparison
proof, so no assertion for an arbitrary fixed ordinary modulus
is implicit. Fix odd $e$ and a smooth sharp $D\in|L^e|$.
For odd $a$, the finite intersection algebra on $P_a=|L^a|$
has determinant line
\[
 \det\pi_*\mathcal O_{C\cap D}
       =\mathcal O_{P_a}(-18ae-9e^2).
\]
The cohomology sequence, removal of the codimension-two
containment locus, and characteristic-two local matrix
$\left(\begin{smallmatrix}1&c\\0&b'\end{smallmatrix}\right)$
are unchanged from Lemma~\ref{lem:ss-intersection-Hasse-degree}.
Thus the reduced irreducible point-intersection divisor
$R_D$ has degree $18ae+9e^2$. For $a\ge3$ the two factor
line bundles separate joint first jets and second jets.
For $a=1$, the product-cubic calculation still gives full
joint curve jets away from finitely many points, rank at
least three at those points, and full second jets at a
nonflex point. These statements follow by tensoring the two
elliptic evaluation spaces; the ordinary cubic is no more
exceptional than the supersingular one. The generic point
of $R_D$ therefore has precisely one contact of order two.
The marked incidence can be generically inseparable; the
local finite-algebra determinant proves the reduced equation
on the coefficient space itself.

For $a\ge3$ the preceding odd lemma gives an irreducible
reduced source Hasse polynomial $h_a$ of degree $9a^2$.
The singular-curve discriminant is irreducible of degree
$c_2(J^1L^a)=54a^2$, by the same first/second-jet calculation
and a generic unique ordinary node. The three degrees
\[
 9a^2,\quad54a^2,\quad18ae+9e^2
\]
are pairwise distinct: either equality involving the last
would require $(a+e)^2=2a^2$ or $7a^2$.
Thus the generic curve on $h_a$ is smooth and transverse
to $D$, and the generic curve on $R_D$ is smooth and sharp.
The respective union defects are at least $2+1=3$ and
exactly $1+1+1=3$. These exceed the generic even defect two.
There are also transverse nodes, which place these unions
outside $Z$ by the mixed-Hessian calculation above.
Consequently both $h_a$ and $r_D$ divide the nonzero restriction
of $h_{a+e}^{\mathrm u}$ under multiplication by $s_D$.
Their degrees add to $9(a+e)^2$, so
\[
                  j_D^*h_{a+e}^{\mathrm u}=c_D h_a r_D,
                  \qquad c_D\ne0.
\]
This uses stable-rank divisibility and total degree, not an
unproved exactness assertion for coherent first-image lattices.

For $a=1,e=3$, the restriction is still nonzero and divisible
by the irreducible $r_D$ of degree $135$. Its generic source
curve is sharp because its Hasse polynomial has degree only
nine, and smooth because the full singular-curve discriminant
has degree $54$. Hence this divisibility needs no irreducibility
of the degree-one source Hasse divisor.

The boundary pairs $(a,e)=(d-1,1)$ and $(d-3,3)$ for $d\ge6$
would force any two global components to have both unordered
degree pairs
\[
 \{9(d-1)^2,18d-9\},\qquad
 \{9(d-3)^2,54d-81\}.
\]
They cannot agree: the crossed equality would give
$(d-4)^2=6$. The first squarefree restriction already forces
the global form to be reduced with at most two components.
For $d=4$, their hypothetical degrees would be $81,63$,
contradicted by the irreducible degree-$135$ restriction
factor just obtained. This proves reduced irreducibility in
every claimed degree. On the generic $R_D$ boundary with
$a=d-1,e=1$ the stable loss is exactly one and the point lies
outside $Z$. Irreducibility and lower semicontinuity of the
rank of a high Frobenius iterate give generic loss one on
the global determinant divisor.
\end{proof}

\begin{lemma}[Two exact mixed degree-two coefficient points]
\label{lem:mixed-degree-two-certificate-candidate}
On the coefficient space for
\[
 E_o:y^2+xy=x^3+1,\qquad E_0:y^2+y=x^3,
 \qquad N=\mathcal O(6O)\boxtimes\mathcal O(6O)
\]
over $\mathbf F_2$, there are two points of the rank-$34$
unit-image open with the following properties. At the first,
the Frobenius on the common kernel of the two invariant
covectors has irreducible nonreciprocal characteristic
polynomial of degree $32$. At the second, every fixed vector
lies in that common kernel.
\end{lemma}

\begin{proof}
Order the factor basis as $(1,x,x^2,x^3,y,xy)$ and its tensor
basis by $6i+j$, with $0\le i,j<6$. A nonnegative integer
encodes its coefficient vector by binary digits in this order.
The section integers are
\[
                 s_A=34346068584,\qquad s_B=68001692233.
\]
Let $C_s$ be the root-removed Cartier multiplication matrix.
The exact independent reconstruction gives
$\rank C_s=\rank C_s^2=34$ at both points. The two fixed
covectors $ns,ens$ have integer vectors
\[
 \begin{array}{c|cc}
  &ns&ens\\\hline
 A&1494603548&420583168\\
 B&26911823620&74729216.
 \end{array}
\]
The full degree-$36$ characteristic polynomial at $A$ is
$T^2(T+1)^2 f(T)$, where
\[
 \begin{split}
 f(T)={}&T^{32}+T^{30}+T^{29}+T^{28}+T^{27}+T^{25}
       +T^{24}+T^{23}+T^{22}\\
       &+T^{20}+T^{18}+T^{15}+T^{14}+T^{11}+T^9+T^8
       +T^5+T+1.
 \end{split}
\]
Its binary polynomial integer is $6372510499$.
The certificate checks $T^{2^{32}}=T\pmod f$ and
$\gcd(T^{2^{16}}-T,f)=1$, proving irreducibility, and checks
that its reciprocal differs from $f$. The prime $65537$
is verified by trial division and has $\operatorname{ord}_{65537}(2)=32$.

At $B$ the literal vectors
\[
                         \mu_0=8601123076,\qquad
                         \mu_1=41642123014
\]
satisfy $(C_{s_B}-I)\mu_0=ns_B$ and
$(C_{s_B}-I)\mu_1=ens_B$.
The coherent Frobenius matrix is $C_s^t$, so these identities
annihilate both covectors on every Frobenius-fixed vector.
Nilpotent directions contain no nonzero fixed vectors;
therefore the assertion holds on the required unit image.
The two covectors are independent, and the rank checks put
both points outside $Z$ and on the unit-image open.
Their curves need not be smooth: restriction from this normal
coefficient open to its dense smooth-member open preserves
the image of every finite cover.

The discovery program and its original output are retained as
\path{agents/monodromy/code/mixed_degree_two_group_discovery.py}
and \path{agents/monodromy/outputs/mixed_degree_two_group_discovery.json}
(306 trials, 51 exact characteristic polynomials,
$4.062721072004933$ seconds).
The independent program
\path{agents/monodromy/code/verify_mixed_degree_two_group.py}
imports no discovery code or stored matrix. It multiplies in
the two coordinate rings and reconstructs Cartier by inverting
the square map on the six-dimensional section space, using
\[
 C_o(f\,dx/x)^2=(1+e)f,\qquad C_0(f\,dx)^2=n f.
\]
It then forms both tensor matrices, computes their integer
characteristic polynomials and reduces modulo two, and checks
all the displayed ranks and vector identities. Its retained
output \path{agents/monodromy/outputs/verify_mixed_degree_two_group.json}
is PASS in $0.4645676829968579$ seconds. These are finite
coefficient checks, not a claim that a single finite-field
fibre has the desired geometric group.
\end{proof}

\begin{lemma}[The geometric generic mixed degree-two image]
\label{lem:generic-mixed-degree-two-group-candidate}
If $j(E)$ is transcendental over $\mathbf F_2$, the degree-two
geometric image is the entire pointwise stabilizer of its
two independent invariant covectors:
\[
                 \operatorname{Hom}(\mathbf F_2^2,\mathbf F_2^{32})
                          \rtimes\GL_{32}(\mathbf F_2).
\]
\end{lemma}

\begin{proof}
Use the family $E_t$ from the first lemma, now restricted to
$t\ne0$, and the maximal-rank even first-image open over it.
It is a normal integral open of a projective bundle. Let $G$
be its full finite arithmetic Lang image, and $H$ the image
of its geometric generic coefficient fibre.
Then $H$ is normal in $G$. Indeed every connected finite
etale cover of the normal total open is integral, and its
generic fibre over the integral curve is a localization of
that integral scheme, hence integral. Thus the generic-fibre
fundamental group surjects onto the total fundamental group.
The geometric generic-fibre fundamental group is normal in
the former, by the exact fundamental-group sequence over
the function field. Passing to finite images proves the claim.

The two covectors are defined over this ordinary base and
fixed arithmetically, including over its imperfect function
field. Explicitly the invariant derivation $D$ on $E_t$ is
\[
 D(x)=tx+1+t,\qquad D(y)=x^2+ty,\qquad D^2=tD.
\]
On $t\ne0$ put $e=D/t$, so $e^2=e$, and use the standard
nilpotent derivation $n=\partial_{x_0}$ on $E_0$.
The invariant volume
\[
             \left(\frac{t\,dx}{tx+1+t}\right)\wedge dx_0
\]
has contraction one with $(e,n)$. The direct local
coefficient-extraction identities of
\eqref{eq:nonordinary-Cartier-table} therefore give the two
fixed covectors $ns,ens$ over this base itself; this does not
require a perfect-field descent argument over $\mathbf F_2(t)$.
At $t=1$ the derivations and volume specialize to exactly
those in the coefficient certificate.
Consequently $G$ fixes the quotient $V/K\simeq\mathbf F_2^2$
pointwise, where $\dim K=32$.
The generic-fibre block lemma gives an actual natural
$\GL_8$ block in $H$, fixing a complement and lying in $K$.
Let $N$ be the normal closure in $G$ of its transvections.
Since $H\triangleleft G$, one has $N\subset H$.
On $K$, the group $N_K$ is nontrivial, normal in $G_K$,
and generated by transvections.

The first certified coefficient point is a rational point of
the total open. Its arithmetic Frobenius, or its inverse if
the opposite convention is used, has the irreducible
nonreciprocal polynomial $f$ or its reciprocal on $K$.
Lemma~\ref{lem:product-frobenius-group-test}, in dimension
$32$, now gives $N_K=G_K=\GL(K)$. In particular $H_K$ is full.

The translation kernel of $G\to\GL(K)$ is a submodule of
$\operatorname{Hom}(\mathbf F_2^2,K)=K\otimes(\mathbf F_2^2)^*$.
Because the natural module has endomorphism ring $\mathbf F_2$,
it has the form $K\otimes U$ for a subspace
$U\subset(\mathbf F_2^2)^*$. The vanishing of the first
cohomology of $\GL(K)$ on sums of its natural module makes
$G$ conjugate to $(K\otimes U)\rtimes\GL(K)$.
If $U$ were proper, this group would have a nonzero fixed
vector outside $K$. The second certified Frobenius element
has no fixed vector outside $K$, excluding that possibility.
Thus $G$ is the entire pointwise stabilizer.

Finally $H$ is normal in this entire stabilizer and surjects
onto $\GL(K)$. Its commutators with all translations include
every element of $K\otimes(\mathbf F_2^2)^*$: for each quotient
coordinate the vectors $(A-I)v$, with $A\in\GL(K)$, span $K$.
Normality puts these commutators in $H$. Hence $H$ contains
all translations and is the entire stabilizer as well.
This step is why the two finite-field coefficient points
give a geometric generic-fibre conclusion, rather than only
a total-family arithmetic conclusion.
\end{proof}

\begin{theorem}[All generic mixed-product images]
\label{thm:generic-mixed-product-all-groups-candidate}
Let $k$ be a perfect field of characteristic two, let $E/k$
be ordinary with $j(E)$ transcendental over $\mathbf F_2$,
and let $X=E\times E_0$ have the Segre embedding defined
by $L$. For every $d\ge1$, the entire geometric and arithmetic
relative Hasse--Witt images are
\[
 G_{\mathrm{geom},d}=G_{\mathrm{arith},d}=
 \begin{cases}
 \GL_{9d^2-1}(\mathbf F_2),&d\text{ odd},\\[2pt]
 \operatorname{Hom}(\mathbf F_2^2,\mathbf F_2^{9d^2-4})
       \rtimes\GL_{9d^2-4}(\mathbf F_2),&d\text{ even}.
 \end{cases}
\]
The even image is the full pointwise stabilizer of its entire
two-dimensional invariant-covector space in rank $9d^2-2$.
Its common covector kernel is irreducible and is contained in
every nonzero invariant subspace; the quotient is the trivial
two-dimensional representation. In particular the odd systems
are geometrically irreducible and the even systems are not.
\end{theorem}

\begin{proof}
Odd degrees and degree two are the preceding lemmas. For
even $d\ge4$, use the reduced irreducible unit-image Hasse
divisor. The determinant-allocation argument of
Lemma~\ref{lem:irreducible-hasse-quotients-trivial} applies
to the negative reflexive sheaf $\mathcal I$ exactly as in
Theorem~\ref{thm:ss-even-full-group-if-unit-Hasse-irreducible}.
In codimension one, determinants of a saturated subobject
and its torsion-free quotient add to the reduced irreducible
determinant divisor. Assigning no divisor to a nonzero
subobject would make its reflexive extension a trivial
Frobenius-unit bundle, contrary to $H^0(\mathcal I)=0$.
Thus every proper nonzero invariant subspace has trivial
quotient. The entire covector space has dimension two, by
Theorem~\ref{thm:nonordinary-ppav-all-covectors}; hence every
such subspace contains its common kernel $K$.
It follows that $K$ is irreducible and that $V$ has no
nonzero fixed vector.

Generic stable loss one gives generation by transvections,
using the Lang finite part and purity argument of
Lemma~\ref{lem:char2-defect-one-inertia}. Removing $Z$ causes
no exception, because its codimension is at least nine.
The restriction to $K$ is thus irreducible and generated
by transvections. Its actual natural $\GL_8$ block excludes
a nondegenerate invariant alternating form: the eight-space
has no invariant form and pairs trivially with its fixed
complement. The exact alternatives of
\cite[Theorem~3.6, printed p.~22]{KantorSeress2001}
therefore give the full $\GL(K)$ image.

As in the degree-two lemma, the translation kernel has the
form $K\otimes U$, and the vanishing of natural-module
$H^1$ gives a linear complement modulo that kernel.
If $U$ were proper, a nonzero quotient direction would
give a fixed vector of $V$. There are no such vectors,
so $U=(\mathbf F_2^2)^*$ and the entire stabilizer occurs.
Both invariant covectors descend over every perfect base
field by the arithmetic part of the cited covector theorem.
Thus the arithmetic image is bounded by the same stabilizer
and contains the geometric image, proving equality.

The plane cubic coordinate rings are generated in degree
one, and the diagonal tensor product has the same property.
Ambient degree-$d$ forms therefore surject onto $H^0(X,L^d)$
for every $d$. The corresponding linear projection has
geometrically integral fibres, and restriction to the dense
smooth maximal-rank open preserves the finite image.
This proves the statement for the requested ambient
hypersurface sections as well as the complete linear systems.
\end{proof}

\begin{remark}[The exact remaining modulus restriction]
\label{rem:generic-mixed-modulus-gate}
The theorem treats every degree at the geometric generic
ordinary modulus, and hence every ordinary elliptic factor
with transcendental $j$. It does not prove these images for
every fixed ordinary elliptic curve over an algebraic extension
of $\mathbf F_2$. Specialization gives the generic conclusion
used above; reversing that direction would be invalid.
The independent exact-rank theorem for every fixed ordinary
factor does not remove this monodromy restriction.
\end{remark}
